\section{Экспериментальное исследование}
\thispagestyle{plain}

% TODO: ~14 страниц --- по результатам реального эксперимента.

\subsection{Методика эксперимента}
% Гипотеза H1: применение разработанного сервиса (с геймификацией
% и интервальным повторением) повышает успеваемость и удержание
% знаний по сравнению с традиционным форматом.
% Дизайн: квазиэкспериментальный с контрольной и экспериментальной группой.
% Этапы: pre-test $\rightarrow$ обучающие сессии $\rightarrow$ post-test
% $\rightarrow$ retention-test через 1--2 недели.
% Этические аспекты: добровольное участие, обезличивание данных.

\subsection{Условия проведения эксперимента}
% Дисциплина: «Надёжность» (преподаватель Волков Д.А., кафедра АСУ).
% Группы: указать конкретные группы и количество участников.
% Период: указать даты pre-/post-/retention-тестов.
% Используемая платформа: разработанный сервис на публичном URL.

\subsection{Метрики}
% Успеваемость: средний балл за тестирование (pre, post, retention).
% Удержание знаний: разница между post и retention.
% Вовлечённость: число активных дней, средняя длительность сессии,
% число добровольных повторений (SM-2).
% Надёжность теста: альфа Кронбаха, дискриминативность вопросов.

\subsection{Статистический анализ}
% Проверка нормальности (Шапиро\textendash Уилка).
% Сравнение средних: t-критерий Стьюдента (если нормальность)
% или критерий Манна\textendash Уитни (если ненормальность).
% Размер эффекта: $d$ Коэна.
% Поправка на множественные сравнения: Бонферрони или FDR.
% Оценка надёжности теста: $\alpha \geq 0{,}7$.

\subsection{Результаты}
% Описательные статистики по группам.
% Графики: распределение баллов, динамика во времени,
% сравнение групп.
% Результаты статистических тестов.
% Расчёт надёжности тестового инструмента.

\subsection{Обсуждение}
% Подтверждение или опровержение гипотезы H1.
% Сопоставление с литературой.
% Ограничения исследования: размер выборки, эффект новизны,
% возможные смещения.
% Практические выводы для образовательного процесса.

\textbf{Выводы по главе 5.}
% Что показал эксперимент, как это согласуется с моделями главы 2.

\newpage
